\section{Conclusion}
\label{sec:conclusion}

We presented \MethodName{}, a method for fast online RL on top of representations extracted from a large pretrained VLA. By training the VLA to expose a compact representation, our method enables a lightweight actor and critic to improve highly precise and delicate tasks with just a few hours of real-world practice. Across four difficult tasks that require precision and speed, \MethodName{} consistently improves both the success rate and execution speed, achieving up to 3× speedup on the hardest phase of each task and, in some cases, surpassing expert human teleoperation speed through strategies that emerge from online RL.

%Ablations confirmed that each component — the RL token, action chunking, reference-action conditioning, and policy regularization — contributes meaningfully, and comparisons against alternative approaches showed that RL methods lacking these elements could fail under long-horizon sparse rewards or produce substantially weaker gains.

While \MethodName{} provides fast and efficient learning, it does require additional human intervention during training to provide reward signals, intervention corrections, and switching between RL (for the critical phase) and the base policy (for the other phases). In principle, some of these components can be automated, for example by using reward models and progress prediction. Developing a fully autonomous RL improvement pipeline based on \MethodName{} is a promising direction for future work. More broadly, we believe that our method represents an important step toward robotic systems that not only learn from demonstration data, but can improve directly on the job. When improvement is fast and reliable, it is enough for the pretraining phase of the VLA to simply provide a good initialization for downstream exploration, while the most successful and performant strategy can be discovered through reinforcement learning. We hope that \MethodName{} will serve as a step toward this future.

%\textbf{Limitations and future work.} \MethodName{} currently relies on a human operator to provide the success/failure labels and trigger the handoff from base policy to RL and to; automating this would make the system easier to deploy for continuous setting. The RL Token is trained so far on task-specific demonstrations, and it remains a question how to transfer it across tasks and robot embodiments. Finally, while \MethodName{} avoids updating the VLA during online RL, finding a way to leverage RL-collected experience to make the foundation model better — closing the loop between fast online adaptation and large-scale model improvement — can be a natural and promising next step